\section{Polynomial factorization: commutatively non-zero case}\label{cnz-sec}

Recall that $\fR$ denotes the free noncommutative polynomial ring
$\F\angle{x_1,x_2,\ldots,x_n}$ and our goal is to give a randomized
polynomial-time factorization algorithm for input polynomials in $\fR$
given as arithmetic formulas when $\F=\F_q$ is a finite field of size
$q$.

A polynomial $f\in\fR$ is \emph{commutatively nonzero} if
$f(\alpha_1,\alpha_2,\ldots,\alpha_n)\ne 0$ for scalars $\alpha_i\in
\F$ (or a small extension field of $\F$).

In this section we will present the factorization algorithm for
commutatively nonzero polynomials.\footnote{In the next section we
  will deal with the general case. The algorithm is more technical in
  detail, although in essence the same.}  It has three broad steps:
\begin{itemize}
\item[(i)] We transform the given polynomial $f$ to a full and \emph{right (or left)
  monic} linear matrix $L$ by first the Higman linearization of $f$
  followed by the algorithm in the proof of Theorem~\ref{full-monic}.
\item[(ii)] Next, we factorize the full and right (or left) monic linear matrix $L$
  into atoms.
\item[(iii)] Finally, we recover the irreducible factors of $f$ from
  the atomic factors of $L$.
\end{itemize}  

We formally state the three problems of interest in this paper.

\begin{problem}[$\prob{FACT}(\F)$]\hfill{~}\\
\textbf{Input}~A noncommutative polynomial $f\in\FX$ given by an
arithmetic formula.\\
\textbf{Output}~Compute a factorization
$f=f_1f_2\cdots f_r$, where each $f_i$ is irreducible, and each $f_i$
is output as an algebraic branching program.
\end{problem}

\begin{problem}[$\prob{LIN{-}FACT}(\F)$]\hfill{~}\\
\textbf{Input}~A full and right (or left) monic linear matrix $L\in\fR^{d\times 
  d}$.\\  
\textbf{Output}~Compute a factorization
$L=F_1F_2\cdots F_r$, where each $F_i$ is a full linear matrix that is
an atom.
\end{problem}

\begin{problem}[$\prob{INV}(\F)$]\hfill{~}\\
\textbf{Input}~A list of scalar matrices $A_1,A_2,\ldots,A_n\in
\F^{d\times d}$.\\
\textbf{Output}~Compute a nontrivial invariant
subspace $V\subset \F^d$ or report that the only invariant subspaces
are $0$ and $\F^d$.
\end{problem}  

In the three-step outline of the algorithm, for the second step we
will show that factoring a full and right (or left) 
monic linear matrix is randomized polynomial-time reducible to the
problem of computing a common invariant subspace for a collection of
scalar matrices. For the third step, we will give a polynomial-time
algorithm (based on Lemma~\ref{triv-lemma}) to recover the irreducible
factors of $f$ from the atomic factors of $L$.

\begin{remark}
  We use Ronyai's randomized polynomial-time algorithm \cite{Ronyai2}
  to solve the problem of computing a a common invariant subspace for
  a collection of matrices over $\F_q$.  Over rational numbers $\Q$,
  even for a special case the problem of computing a common invariant
  subspace turns out to be at least as hard as factoring square-free
  integers~\cite{FR85}. Hence, our approach to noncommutative
  polynomial factorization does not yield an efficient algorithm over
  $\Q$.
\end{remark}

Suppose $f\in\fR$ is given by a noncommutative arithmetic
formula. Since $f$ has small degree we can check if it is
commutatively nonzero in randomized polynomial-time by the
DeMillo-Lipton-Schwatrtz-Zippel test \cite{DL78,Sch80,Zip79} and, if
so, find $\alpha_i\in\F, i\in[n]$ such that
$f(\alpha_1,\alpha_2,\ldots,\alpha_n)\ne 0$ (if $\F$ is small, we pick
$\alpha_i$ from a small extension field). Furthermore, by a linear
shift of the variables $x_i\leftarrow x_i + \alpha_i, i\in[n]$
followed by scaling we can assume $f(\overline{0}) = 1$. Note that
from the factorization of the linear shift of $f$ we can recover the
factors of $f$ by shifting the variables back, and irreducibility is
preserved by linear shift. For the rest of this section we will
assume $f(\overline{0})=1$. 

Let $L=A_0+\sum_{i=1}^n A_i x_i$. As $f(\overline{0})=1$, we have
$L(\overline{0})=A_0$ is an invertible matrix. We now present an
efficient algorithm for factoring $L$ as a product of linear matrices
$L_1L_2\cdots L_r$, where each $L_i$ is an atom.

\begin{remark}
The factorization algorithm for arbitrary full and right (or left)
monic linear matrices (in which $A_0$ need not be invertible) is
similar but more involved. It is based on Lemma~\ref{shift-inv-lemma}
and is dealt with in the next section.
\end{remark}

\subsection{Algorithm for a special case of $\prob{LIN{-}FACT}(\F_q)$}

%
%
%
%
%

\begin{theorem}\label{lfact1}
  There is a randomized polynomial-time algorithm for the following
  two special cases of the $\prob{LIN{-}FACT}(\F_q)$ problem:
  \begin{enumerate}
 \item Given a full right monic matrix $L$ as input such that
   $L(\overline{0})$ is an invertible matrix, the algorithm outputs a
   factorization of $L$ as a product of linear matrices that are
   atoms.
  \item Given a full left monic matrix $L$ as input such that
   $L(\overline{0})$ is an invertible matrix, the algorithm outputs a
   factorization of $L$ as a product of linear matrices that are
   atoms.  
\end{enumerate}
\end{theorem}  

\begin{proof}
  We present the algorithm only for the first part, as the second part
  has essentially the same solution.
  
  Let $L=A_0+\sum_{i=1}^n A_i x_i$ in $\fR^{d\times d}$ be such an
  instance of $\prob{LIN{-}FACT}(\F_q)$. We can write $L=A_0\cdot
  L'$ where $L'$ is the full and right monic linear matrix
  \[
  L'= I_d + \sum_{i=1}^n A_0^{-1}A_i x_i.
  \]
 Clearly, it suffices to factorize the linear matrix $L'$ into
 atoms. 
 
First we show that $L'$ is an not atom iff matrices $A_0^{-1}A_i$, $1\leq i \leq n$ have a nontrivial common invariant subspace.  
 By Theorem~\ref{cohnthm}, $L'$ is not an atom if and only if
 we can write $S_1L'S_2 = \left(
\begin{array}{cc}
B & 0 \\
D & C
\end{array}
\right)$ for invertible scalar matrices $S_1$ and $S_2$, where $B$ and
$C$ are full and right monic linear matrices, and $D$ is some linear
matrix. Equating the constant terms on both sides of 
the above equation we have $S_1S_2= \left(
\begin{array}{cc}
  B_0 & 0 \\
  D_0 & C_0
\end{array}
\right)$ as the constant term of $L'$ is $I_d$.  Thus the matrices $S_1S_2$ and its inverse also has the
same block form which implies that
$S_1L'S_1^{-1}=S_1L'S_2(S_1S_2)^{-1}$ also has the same block form. It follows that the $n$ matrices $A_0^{-1}A_i, 1\le i\le n$ have
a nontrivial common invariant subspace. Conversely, if the matrices
$A_0^{-1}A_i, 1\le i\le n$ have a nontrivial common invariant subspace
then we have a basic change scalar matrix $S$ such that $SL'S^{-1}$ has
the block form $\left(
\begin{array}{cc}
  L_1 & 0 \\
  * & L_2
\end{array}
\right)$,
where $L_1$ and $L_2$ are full and right monic linear matrices. So by Theorem~\ref{cohnthm} $L'$ is not an atom. 
 So we have established, $L'$  (and hence $L$) is not an atom iff matrices $A_0^{-1}A_i$, $1\leq i \leq n$ have a nontrivial common invariant subspace.
We will use Ronyai's randomized polynomial-time algorithm for finding
a nontrivial common invariant subspace for matrices $A_0^{-1}A_i, 1\le
i\le n$ over finite field $\F_q$. 

If there is no nontrivial invariant
subspace then the linear matrix $L'$ (and hence $L$) is an
atom. 
Otherwise, by repeated application of Ronyai's algorithm we will
obtain a basis change scalar matrix $T$ which when applied to $L'$
yields a linear matrix in the following \emph{atomic block diagonal
  form}:
\begin{equation}\label{eq3}
TL' T^{-1} = 
\left(
\begin{array}{ccccc}
L_1 & 0 &0 &\ldots & 0\\
* & L_2 &0 &\ldots & 0\\
* & * &L_3 &\ldots & 0\\
& & &\ddots &\\
* & * &* &\ldots & L_r\\
\end{array}
\right),
\end{equation}
where for each $j\in[r]$, the full right monic linear matrix $L_j \in \fR^{d_j \times
  d_j}$ is an atom, and each $*$ stands for some unspecified linear
matrix. It is now easy to factorize $TL'T^{-1}$ as a product of atoms
by noting one step of the factorization of $TL'T^{-1}$ from its form:
\[
TL'T^{-1}=  \left(
\begin{array}{cc}
A & 0 \\
D & L_r
\end{array}
\right) =   \left(
\begin{array}{cc}
A & 0 \\
0 & I
\end{array}
\right)\cdot 
  \left(
\begin{array}{cc}
I & 0 \\
D & I
\end{array}
\right)\cdot
\left(
\begin{array}{cc}
I & 0 \\
0 & L_r
\end{array}
\right).
\]
We note that  $\left(
\begin{array}{cc}
I & 0 \\
D & I
\end{array}
\right)$ is a unit. Since $L_r$ is an atom the product $\left(
\begin{array}{cc}
I & 0 \\
D & I
\end{array}
\right)\cdot \left(
\begin{array}{cc}
I & 0 \\
0 & L_r
\end{array}
\right)$ is also an atom and a linear matrix, and it is the rightmost
factor of $TL'T^{-1}$. Continuing thus with $A$ now, we can factorize
$TL'T^{-1}$ as a product $F'_1F'_2\cdots F'_r$ of $r$ atoms, each of
which is a linear matrix. It follows that $L=A_0T^{-1}F'_1F'_2\cdots
F'_rT$ is a complete factorization of $L$ as a product of atomic
linear matrices (both $A_0$ and $T$ are scalar invertible matrices).
\end{proof}

\begin{remark}
We note that Ronyai's algorithm \cite{Ronyai2} for $\prob{INV}(\F_q)$
is actually a deterministic polynomial-time \emph{reduction} from
$\prob{INV}(\F_q)$ to univariate polynomial factorization over $\F_q$.
\end{remark}  

 
%
%
%

Based on whether we want to work with right monic or left
monic case we will express $f \oplus I_s$ in an appropriate form using
Higman linearization and Theorem~\ref{full-monic} as described in the
equation below:
\begin{equation}\label{eq2}
f\oplus I_s = 
\begin{cases}
~~PU(L'\oplus I_t) SQ,~ \text{in the right monic case}\\ 
~~PS(L'\oplus I_{t}) UQ,~ \text{in the left monic case} 
\end{cases}
\end{equation}


where $d+t=s+1$, $L'\in\fR^{d\times d}$ is a full and right (or left) monic linear
matrix, $P$ is upper triangular with all $1$'s diagonal, $Q$ is lower
triangular with all $1$'s diagonal, $U\in\fR^{(d+t)\times (d+t)}$ is a
unit, and $S\in\F^{(d+t)\times (d+t)}$ is an invertible scalar
matrix. 


\subsection*{Algorithm for $\prob{FACT}(\F_q)$}

We are now ready to describe the polynomial factorization algorithm
for commutatively nonzero polynomials in $\fR$. Starting with the
Higman linearization of the input polynomial $f\in\fR$ as in
Equation~\ref{eq2}, by an application of the first parts of Theorems
\ref{full-monic} and \ref{lfact1} we obtain the factorization
$f\oplus I_s = PUF'_1F'_2\cdots F'_r SQ$ using the structure in
Equation~\ref{eq3}. 

Alternatively, by applying the second part of Theorem \ref{full-monic}
we can compute a left monic linear matrix $L'$ that is a stable
associate of $f$ and, applying the second part of Theorem~\ref{lfact1}
we can compute the factorization
\begin{equation}\label{eq4}
f\oplus I_s = PS'F'_1F'_2\cdots F'_r U'Q.
\end{equation}
where each linear matrix $F'_i$ is an atom, $P$ is upper triangular
with all $1$'s diagonal, $Q$ is lower triangular with all $1$'s
diagonal, $U'$ is a unit and $S'$ is a scalar invertible matrix.
Equation~\ref{eq4} is the form we will use for the algorithm (we could
equally well use the other factorization).

{From} the structure of the atomic block diagonal matrix $TL'T^{-1}$
in Equation~\ref{eq3} notice that the product $S'F'_1F'_2\cdots F'_i $
is a linear matrix for each $1\le i<r$.

The next lemma presents an algorithm that is crucial for extracting
the factors of $f$.

\begin{lemma}\label{triv-lem}
  Let $C\in\fR^{u\times d}$ be a linear matrix and $v\in\fR^{d\times
    1}$ be a column of polynomials such that $Cv=0$. Each entry $v_i$ of $v$ is given by an
  algebraic branching program as input. Then, in polynomial time we
  can compute a invertible matrix $N\in\fR^{d\times d}$ such that
  \begin{itemize}
  \item For $1\le i\le d$ either the $i^{th}$ column of $CN$ is all zeros
    or the $i^{th}$ row of $N^{-1}v$ is zero.
  \item Each entry of $N$ is a polynomial of degree at most $d^2$ and
    is computed by a polynomial size ABP, and also each entry of
    $N^{-1}$ is computed by a polynomial size ABP.
  \end{itemize}
%
\end{lemma}

\begin{proof}
 We will describe the algorithm as a recursive procedure Trivialize
 that takes matrix $C$ and column vector $v$ as parameters and returns
 a matrix $N$ as claimed in the statement.

  \begin{enumerate}
  \item[] Procedure Trivialize$(C\in\fR^{u\times d},v\in\fR^{d\times 1})$
  \item If $d=1$ then (since $Cv=0$ iff either $C=0$ or $v=0$)
    \textbf{return} the identity matrix.
  \item  If $d>1$ then
  \item write $C=C_0+C_1$, where $C_0$ is a scalar matrix and $C_1$ is the
    degree $1$ homogeneous part of $C$. Let $k$ be the degree of the
    highest degree nonzero monomials in the polynomial vector $v$, and
    let $m$ be a nonzero degree $k$ monomial. Let
    $v(m)\in\F_q^{d\times 1}$ denote its (nonzero) coefficient in $v$.
    Then $Cv=0$ imples $C_1 v(m)=0$. Let $T_0\in\F_q^{d\times d}$ be a
    scalar invertible matrix with first column $v(m)$ obtained by
    completing the basis.
    \begin{enumerate}
    \item If $C_0v(m)=0$ then the first column of $CT_0$ is zero.
    \item Otherwise, $CT_0$ has first column as the nonzero scalar
      vector $Cv(m)=C_0v(m)$. Suppose $i^{th}$ entry of $Cv(m)$ is a
      nonzero scalar $\alpha$. With column operations we can drive the
      $i^{th}$ entry in all other columns of $CT_0$ to zero. Let the
      resulting matrix be $CT_0T_1$ (where the matrix $T_1$ is
      invertible as it is a product of elementary matrices
      corresponding to these column operations, each of which is of
      the form $\col_i\leftarrow (\col_i +
      \col_1\cdot \alpha_0+\sum_i \alpha_i x_i)$). Notice that $CT_0T_1$ is still
      linear.
    \item As $Cv=(CT_0T_1)(T_1^{-1}T_0^{-1}v)$, and in the $i^{th}$
      row of $CT_0T_1$ the only nonzero entry is $\alpha$ which is in
      its first column, we have that the first entry of
      $T_1^{-1}T_0^{-1}v$ is zero.
    \end{enumerate}
  \item Let $C'\in\fR^{u\times (d-1)}$ obtained by dropping the first
    column of $CT_0T_1$. Let $v'\in\fR^{(d-1)\times 1}$ be obtained by
    dropping the first entry of $T_1^{-1}T_0^{-1}v$. Note that $C'$
    is still linear.
  \item Recursively call Trivialize$(C'\in\fR^{u\times
    (d-1)},v'\in\fR^{(d-1)\times 1})$.  and let the matrix returned by
    the call be $T_2\in\fR^{(d-1)\times (d-1)}$.
  \item Putting it together, return the matrix $T_0T_1(I_1\oplus T_2)$.
  \end{enumerate}

  To complete the proof, we note that a highest degree monomial $m$
  such that $v(m)\ne 0$ is easy to compute in deterministic polynomial
  time if each $v_i$ is given by an algebraic branching program using
  the PIT algorithm of Raz and Shpilka \cite{RS05}. Notice that for
  the recursive call we need $C'$ to be also a linear matrix and each
  entry of $v'$ to have a small ABP. $C'$ is linear because $CT_0T_1$
  is a linear matrix since $CT_0$ is linear, its first column is
  scalar, and each column operation performed by $T_1$ is scaling the
  first column of $CT_0$ by a linear form and subtracting from another
  column of $CT_0$. Each entry of $v'$ has a small ABP because
  $T_0^{-1}$ is scalar and it is easy to see that the entries of
  $T_1^{-1}$ have ABPs of polynomial size. Finally, we note that $T_1$
  is a product of at most $d-1$ linear matrices (each corresponding to
  a column operation), and $N$ is an iterated product of $d$ such
  matrices. Hence, each entry of $N$ as well as $N^{-1}$ is a
  polynomial of degree at most $d^2$ and is computable by a small ABP.


  
\end{proof}

%
%
%
%

Turning back to our algorithm for $\prob{FACT}(\F_q)$, in the next
lemma we design an efficient algorithm that will allow us to extract
all the irreducible factors of $f$ (given Equation~\ref{eq4}).

\begin{lemma}[Factor Extraction]\label{extract}
  Let $f\in\fR$ be a polynomial and $G\in\fR^{(d-1)\times (d-1)}$ be a
  unit such that
  \begin{equation}\label{eq5}
  \left(
\begin{array}{cc}
f & u \\
0 & G
\end{array}
\right)
= P CD,
\end{equation}
such that
\begin{itemize}
\item $C$ is a full linear matrix that is a non-unit, $P$ is upper
  triangular with all $1$'s diagonal, and $D\in\fR^{d\times d}$ is a
  full non-unit matrix which is also an atom.
\item The polynomial $f$, and the entries of $u, G, P, D$ are all
  given as input by algebraic branching programs.
\end{itemize}
Then we can compute in deterministic polynomial time a nontrivial
factorization $f=g\cdot h$ of the polynomial $f$ such that $h$ is an
irreducible polynomial.
\end{lemma}

\begin{proof}
Let
\[C=\left(
\begin{array}{cc}
c_1 & c_3 \\
c_2 & c_4
\end{array}
\right) \text{ and }
D = 
  \left(
\begin{array}{cc}
d_1 & d_3\\
d_2 & d_4
\end{array}
\right),
\]
written as $2\times 2$ block matrices where $c_1$ and $d_1$ are
$1\times 1$ blocks. By dropping the first row of the matrix in the
left hand side of Equation~\ref{eq5} and the first row of $P$ we
get
\[
  (0~G) = (0~P') C D,
\]
where $P'$ is also an upper triangular matrix with all $1$'s diagonal.
Equating the first columns on both sides we have
\begin{eqnarray*}
  0 & =&  (0~P')\left(
\begin{array}{cc}
c_1 & c_3 \\
c_2 & c_4
\end{array}
\right)\left(\begin{array}{c} d_1\\ d_2 \end{array} \right),
  \text{ which implies that}\\
  0 & = & P'(c_2~c_4)\left(\begin{array}{c} d_1\\ d_2 \end{array} \right),
  \text{ and hence}\\
  0 & = & (c_2~c_4)\left(\begin{array}{c} d_1\\ d_2 \end{array} \right),
  \text{ since $P'$ is invertible.}
\end{eqnarray*}

%

Since $(c_2~c_4)\in\fR^{(d-1)\times d}$ is a matrix with linear
entries and $\left(\begin{array}{c} d_1\\ d_2 \end{array} \right)\in\fR^{d\times 1}$ is a column vector of
polynomials which are given by ABPs as input, we can apply the
algorithm of Lemma~\ref{triv-lem} to compute a unit
$N$ such that its entries are all given by
ABPs such that for $1\le i\le d$, either the $i^{th}$ column of
$(c'_2~c'_4)=(c_2~c_4)N$ is zero or the $i^{th}$ row of
$\left(\begin{array}{c} d_1'\\ d_2' \end{array}
 \right)=N^{-1}\left(\begin{array}{c} d_1\\ d_2 \end{array}
 \right)$ is zero.

%
%
%

Now the following argument is almost identical with the argument towards the end of the proof of the Theorem \ref{associates-preseve-atoms}. We give it below for completeness. 
Since $D$ is a full matrix, the matrix $N^{-1}D$ is also full which
implies its first column $\left(\begin{array}{c} d_1'\\ d_2' \end{array} \right)$ cannot be all zeros. So there is at least one nonzero entry in $\left(\begin{array}{c} d_1'\\ d_2' \end{array} \right)$ and the corresponding column in $(c'_2~c'_4)$ is all zero. This implies there exist a permutation matrix $\Pi$ such that the first column of $C(c'_2~c'_4)\Pi$ is all zero and first entry of $\Pi^{-1} \left(\begin{array}{c} d_1'\\ d_2' \end{array} \right)$ is non zero.

%
%
%
%
%
%
%
%
%
%
%
%
%

Consider the matrices $C''=CN\Pi=\left(
\begin{array}{cc}
c''_1 & c''_3 \\
c_2'' & c''_4
\end{array}
\right)
$ and
$D''=\Pi^{-1}N^{-1}D=\left(
\begin{array}{cc}
d''_1 & d''_3 \\
d_2'' & d''_4
\end{array}
\right)
$. We have

\[  \left(
\begin{array}{cc}
f & * \\
0 & G'
\end{array}
\right)
=P^{-1}
\left(
\begin{array}{cc}
f & u \\
0 & G
\end{array}
\right)
=
\left(
\begin{array}{cc}
c''_1 & c''_3 \\
c_2'' & c''_4
\end{array}
\right)
  \left(
\begin{array}{cc}
d''_1 & d''_3\\
d_2'' & d''_4
\end{array}
\right)
\]

, where $G'=(P')^{-1}G$ is a unit, $c_2''$ is all zero column matrix and $d_1''$ is non-zero. Now observing $(2,1)^{th}$ matrix block in the above equation, we get $d_2''$ is all zero column. Hence, by looking at $(2,2)^{th}$ block in the above equation, we can see that $c_4''$ and $d_4''$ are units as $G'$ is a unit. Clearly, we have $f= c_1'' \cdot d_1 ''$. Now,
since $C$ and $D$ are non-units (by assumption), the matrices $C''$
and $D''$ are also non-units. Therefore, $c''_1$ is not a scalar for
otherwise $C''$ would be a unit. Similarly, $d''_1$ is not a scalar.
It follows that $f=c''_1d''_1$ is a nontrivial factorization of $f$.

Furthermore, since $D$ is an atom by assumption and $D''$ is a stable
associate of $D$, $D''$ is an atom. As $D'' = \left(
\begin{array}{cc}
d''_1 & d''_3\\
0 & d''_4
\end{array}
\right)$ and $d_4''$ is invertible, we get $\left(
\begin{array}{cc}
1 & 0\\
0 & (d''_4)^{-1}
\end{array}
\right) \cdot D'' = \left(
\begin{array}{cc}
d''_1 & d''_3\\
0 & I_{s}
\end{array}
\right)$. Now applying suitable row operations to the matrix $(1 \oplus (d''_4)^{-1})D''$ we can drive $d_3''$ to zero. So we have $U(1 \oplus (d''_4)^{-1})D''=(d_1'' \oplus I_s)$ for a unit $U$. Hence $d_1''$ is an associate of $D''$ and therefore $d_1''$ is irreducible as $D''$ is an atom. 
  

%
%
%
%
%
%
%
%
%
%
%
%
%
%
%
%
%
%
%
%
%
%
%
%
%
%
%
%
%
%
%
%
%
%

%
%
%
\end{proof}

Finally, we describe the factorization algorithm for commutatively
nonzero polynomials $f\in\fR$ over finite fields $\F_q$.

\begin{theorem}\label{cnzthm}  
  Let $\fR=\F_q\angle{X}$ and $f\in \fR$ be a commutatively nonzero
  polynomial given by an algebraic branching program of size $s$ as
  input instance of $\prob{FACT}(\F_q)$. Then there is a
  $\poly(s, \log q)$ time randomized algorithm that outputs a
  factorization $f=f_1f_2\cdots f_r$ such that each $f_i$ is
  irreducible and is output as an algebraic branching program.
\end{theorem} 

\begin{proof}
  Given $f$ as input, we apply Higman linearization followed by the
  algorithm for $\prob{LIN{-}FACT}(\F_q)$ described in
  Theorem~\ref{lfact1} to obtain the factorization of $f\oplus I_s=PSS_1F_1F_2 \ldots F_rS_2UQ$ where 
each linear matrix $F_i$ is an atom, $P$ is upper triangular
with all $1$'s diagonal, $Q$ is lower triangular with all $1$'s
diagonal, $U$ is a unit and $S$ is a scalar invertible matrix, as given in Equation~\ref{eq4}. 
%
   We can now apply Lemma~\ref{extract} to
  extract irreducible factors of $f$ (one by one from the right).

  For the first step, let $C=S S_1F_1F_2\cdots F_{r-1}$ and $D=F_rS_2UQ$
  in Lemma~\ref{extract}. The proof of Lemma~\ref{extract} yields the
  matrix $N_r=N\Pi$ such that both matrices $C''=PS S_1F_1F_2\cdots
  F_{r-1}N_r$ and $D''=N_r^{-1}F_rS_2UQ$ has the first column all zeros
  except the $(1,1)^{th}$ entries $c''_1$ and $d''_1$ which yields the
  nontrivial factorization $f=c''_1d''_1$, where $d''_1=f_r$ is 
  irreducible. Renaming $c''_1$ as $g_r$ we have from the structure of
  $C''$:
  \[
  \left( \begin{array}{cc} g_r & * \\ 0 & G_r \end{array} \right) =
  P(SS_1F_1F_2\cdots F_{r-2}) (F_{r-1}N_r).
  \]
Setting $C= SS_1F_1F_2\cdots F_{r-2}$ and $D=F_{r-1}N_r$ in
Lemma~\ref{extract} we can compute the matrix $N_{r-1}$ using which we
will obtain the next factorization $g_r=g_{r-1}f_{r-1}$, where
$f_{r-1}$ is irreducible because the linear matrix $F_{r-1}$ is an
atom. Lemma~\ref{extract} is applicable as all conditions are met by
the matrices in the above equation (note that $G_r$ will be a unit).

Continuing thus, at the $i^{th}$ stage we will have
$f=g_{r-i+1}f_{r-i+1}f_{r-i+2}\cdots f_r$ after obtaining the
rightmost $i$ irreducible factors by the above process. At this stage
we will have
  \[
  \left( \begin{array}{cc} g_{r-i+1} & * \\ 0 & G_{r-i+1} \end{array}
  \right) = P(SS_1F_1F_2\cdots F_{r-i-1}) (F_{r-i}N_{r-i+1}),
  \]
  where $G_{r-i+1}$ is a unit and all other conditions are met to
  apply Lemma~\ref{extract}.

  Thus, after $r$ stages we will obtain the complete factorization
  $f=f_1f_2\cdots f_r$. For the running time, it suffices to note that
  the matrix $N$ computed in Lemma~\ref{extract} is a product of
  degree at most $d^2$ many linear matrices (corresponding to the
  column operations). Thus, at the $i^{th}$ of the above iteration,
  the sizes of the ABPs for the entries of $N_{r-i+1}$ are independent
  of the stages. Hence, the overall running time is easily seen to be
  polynomial in $s$ and $\log q$.
\end{proof}

\begin{corollary}\label{sparse-cor1}
If $f\in\fR$ is commutatively nonzero polynomial given as input in
sparse representation (as an $\F_q$-linear combination of its
monomials) then in randomized polynomial time we can compute
a factorization into irreducible factors in sparse representation.
\end{corollary}

\begin{proof}
  Let $f$ be given as input in sparse representation. Suppos $\deg f =
  d$ and it is $t$-sparse. Then there are at most $td^2$ many
  monomials that can occur as a substring of the monomials of $f$.  We
  can apply the randomized algorithm of Theorem~\ref{cnzthm} to obtain
  the factorization $f=f_1f_2\cdots f_r$, where each $f_i$ is given by
  an ABP. Now, for each of the $td^2$ many candidate monomials of
  $f_i$ we can find its coefficient in $f_i$ in polynomial time (using
  the Raz-Shpilka algorithm \cite{RS05}). Hence we can obtain the
  factorization $f=f'_1f'_2\cdots f'_r$, where each $f'_i$ is a
  $t$-sparse polynomial.
\end{proof}
